\section{Antecedentes}
\label{sec:antecedentes}
La formulación de políticas públicas contemporáneas opera bajo el mandato global de \textit{``No dejar a nadie atrás''} (\textit{Leave No One Behind}, LNOB), eje transversal de la Agenda 2030 para el Desarrollo Sostenible \cite{cepal2024, un2030}. Este compromiso internacional ha impulsado una transición analítica desde los agregados nacionales hacia una mayor granularidad subnacional, reconociendo que los promedios macroeconómicos suelen disimular disparidades territoriales críticas. En este escenario, la Estimación en Áreas Pequeñas (\textit{Small Area Estimation}, SAE) surge como un paradigma metodológico indispensable para superar las limitaciones de representatividad de las encuestas de hogares, permitiendo fundamentar la intervención estatal en inferencias precisas a nivel municipal y local \cite{cepal2024, dane2024}.

\subsection{Evolución de la medición de pobreza y el desafío de la desagregación territorial}

Los paradigmas de medición de la pobreza han transitado de enfoques estrictamente monetarios hacia marcos multidimensionales basados, principalmente, en la metodología de Alkire y Foster \cite{alkire2011}. Este enfoque constituye la base técnica de diversos índices oficiales y sustenta ejercicios aplicados como el Índice de Pobreza Energética Multidimensional (IMPE) en Colombia, cuya implementación busca trasladar el análisis desde una perspectiva nacional hacia una escala municipal, idónea para identificar brechas estructurales en zonas rurales dispersas \cite{promigas2025}.

No obstante, la generación de estas estadísticas enfrenta un obstáculo inferencial persistente: las encuestas de hogares tradicionales ---como la Pesquisa Nacional por Amostra de Domicílios Contínua (PNADC) en Brasil o la Gran Encuesta Integrada de Hogares (GEIH) en Colombia--- están diseñadas para ofrecer representatividad en dominios geográficos mayores. Por tanto, al desagregar los datos para municipios o estratos territoriales pequeños, la varianza y los errores de muestreo se incrementan sustancialmente, invalidando las estimaciones directas \cite{moura2021, dane2024}.

Según la documentación técnica del Departamento Administrativo Nacional de Estadística (DANE), la arquitectura estadística en Colombia padece una marcada asimetría temporal y geográfica. Mientras el nivel departamental cuenta con una periodicidad anual respaldada por encuestas continuas, la información municipal depende de censos nacionales realizados cada diez o quince años \cite{dane2024}. Este desfase intercensal genera un vacío de información que obstruye el monitoreo oportuno de las metas de desarrollo. Ante esta restricción, la integración de encuestas de alta periodicidad con registros administrativos e información auxiliar se convierte en una condición necesaria para mantener la vigencia de las estadísticas oficiales y garantizar decisiones basadas en evidencia contemporánea.

\subsection{Modelos de estimación en áreas pequeñas en el contexto latinoamericano}

En América Latina, la Comisión Económica para América Latina y el Caribe (CEPAL) ha liderado la transferencia de metodologías SAE, priorizando modelos a nivel de unidad como el \textit{Empirical Best Predictor} (EBP) y el \textit{Pseudo-EBP} \cite{cepal2024}. Este último resulta crucial por su capacidad para incorporar los pesos de muestreo en el modelado, mitigando el sesgo derivado de ignorar el diseño complejo de las encuestas. Estos enfoques articulan la precisión analítica de los datos muestrales con la cobertura exhaustiva de los censos y otras fuentes auxiliares, incluyendo sensores remotos.

Dentro de las experiencias regionales, Brasil constituye un referente metodológico relevante mediante el desarrollo de modelos SAE bajo distribución Beta con función de enlace logística para estimar la incidencia, el hiato y la severidad de la pobreza \cite{moura2021}. Las aplicaciones basadas en la PNADC demuestran que la inclusión de series históricas de al menos cuatro años en estructuras temporales estabiliza y mejora significativamente la precisión frente a los estimadores directos, alcanzando estándares de calidad aptos para la publicación oficial.

A este ecosistema metodológico se suma la incorporación de variables auxiliares geoespaciales. Plataformas como Google Earth Engine facilitan la extracción de predictores robustos ---luces nocturnas, cobertura urbana o índices de vegetación--- a partir de imágenes satelitales como Sentinel-2 y PlanetScope \cite{cepal2024, dane2024}. Estas variables operan como indicadores indirectos (*proxies*) de la actividad económica y el bienestar en territorios con registros administrativos deficientes. Así, la integración de componentes espaciales y fuentes alternas reduce la dependencia de censos desactualizados frente a dinámicas territoriales aceleradas, como la expansión urbana o las transformaciones agrícolas.

\subsection{Avances y experiencias de aplicación en Colombia}

En el ámbito nacional, el DANE ha avanzado en la implementación de modelos de área de tipo Fay-Herriot, aplicados a la Encuesta Nacional de Demografía y Salud (ENDS 2015) para estimar la proporción de hogares con miembros en el extranjero \cite{dane2024}. Asimismo, la entidad evalúa modelos mixtos generalizados con componentes geoestadísticos para el mapeo subnacional de la pobreza monetaria. Aunque estos desarrollos denotan un progreso técnico importante, su consolidación metodológica y difusión pública permanecen en fase de validación.

Por otra parte, el sector privado ha aportado análisis territoriales sobre privaciones multidimensionales. El reporte del IMPE 2025 de Promigas, por ejemplo, cartografía la denominada \textit{periferia energética} en departamentos como Chocó, La Guajira y Vichada \cite{promigas2025}. Los hallazgos revelan que la pobreza energética en zonas rurales remotas supera críticamente a la de los centros urbanos. Frente a este panorama, el informe define metas de reducción del indicador y propone líneas de acción concretas, tales como la optimización del servicio eléctrico, la sustitución de leña, la digitalización del hogar y la electrificación de instituciones educativas.

En el plano tecnológico, el DANE ha incorporado algoritmos de *Machine Learning* como Random Forest y XGBoost para la clasificación supervisada de coberturas vegetales y la actualización de marcos muestrales agropecuarios \cite{dane2024}. Estas innovaciones responden a la exigencia del Plan Nacional de Desarrollo 2022--2026 de producir estadísticas oficiales más desagregadas para visibilizar a grupos históricamente excluidos.

\subsection{Vacíos en la literatura y justificación metodológica}

A pesar de las iniciativas descritas, la estimación de la pobreza multidimensional a nivel municipal mediante modelos mixtos complejos con componentes espaciales sigue siendo un campo en consolidación \cite{dane2024}. Si bien se registran avances en pobreza monetaria y metodologías SAE generales, Colombia aún no dispone de un protocolo estandarizado, validado y publicado para la estimación anualizada de la pobreza multidimensional municipal a partir de las fuentes del DANE.

La literatura identifica, además, desafíos persistentes en la integración de datos. Destacan el sesgo de enlace (*linking bias*), asociado a las dificultades operativas en el cruce de encuestas y registros administrativos, y la cobertura incompleta de estos últimos. Asimismo, se evidencia una ausencia de protocolos consolidados para la homologación de categorías, variables y formatos entre las distintas fuentes de información \cite{dane2024}, lo que compromete la consistencia de los modelos y limita la comparabilidad local.

Ante la inexistencia de una metodología oficial estandarizada para evaluar anualmente el IPM municipal, esta propuesta resulta plenamente justificada. El desarrollo de una estrategia de estimación para áreas pequeñas con componentes espaciales permitirá mitigar el vacío de información intercensal y estructurar un monitoreo sistemático de los riesgos y desigualdades territoriales, proveyendo un insumo robusto para el diseño de políticas públicas focalizadas.